\section{局限性}
\label{sec:limitations}

\subsection{载体参数是合成扩展}

83 个载体工具上的 285 个可选敏感字段参数是本项目新增的，不是四个来源的原生参数。它们解决了参数级标签歧义——每个新增参数精确对应一个档案字段，避免了把多个敏感值拼入一个通用文本字段的问题。但这些参数不能反映生产接口中的自然泄露机会，因为真实 API 的可选参数设计可能与本项目不同。后续应建立原生可选参数子集并进行领域专家审核。

\subsection{参考调用的局限}

参考调用的参数是否为业务上的最小集合，本文不作判断。当前授权规则只根据具体值是否出现在源参考参数中生成标签：如果字段值在参考调用里出现过，就标记为允许。但参考调用本身可能包含多余的参数——上游基准的参考调用并非经过逐项消融证明的最小集合。因此，允许标签的含义是"上游参考使用了该值"，而非"该值是业务上必需的"。

\subsection{环境验证的范围}

328 条发布调用全部完成环境执行，其中 321 条通过最终状态检查。环境验证确认调用可执行且最终状态正确，但不涉及参数是否最小。一个调用可以通过最终状态检查，却仍然携带了不必要的参数——这正是 RPL 的定义。API-Bank 的 7 条因上游不提供整对话最终状态判据而不判定，只能确认每步调用通过上游检查器。

\subsection{档案补全改变上下文}

2{,}280 个档案字段中有 307 个由程序确定性补全，来自 BFCL、API-Bank 与 AppWorld。这些补全值是虚构的，但下游系统看到这些值后可能产生源任务没有的推断。例如，一个原本不涉及医疗信息的任务，补全后档案中可能包含诊断字段，智能体因此可能把诊断信息写入工具参数。实验应同时报告完整档案与 147 条无生成字段子集的结果。

\subsection{任务书链长解释}

正式数据未强制 5/7 链长，而是保留完整参考调用。若只保留自然符合 5/7 的任务，328 条会降至 77 条，AppWorld 仅剩 3 条。固定 5/7 会系统性地删除多步任务，而这些任务恰恰是多工具协作的典型场景。若验收系统硬编码 5/7，应从完整数据导出兼容子集，而不是截断或伪造调用。

\subsection{来源特定偏差}

各来源的数据特性不同，跨来源比较须按来源分层。API-Bank 的可见范围绑定在记录对话范围，对话聚合是本项目的设计而非上游原生功能。AppWorld 的长执行日志主导调用数和链长分布，但注入由显式白名单与每链两步上限控制。BFCL 的工具数量多（129 个），导致禁止关系数量大，但这不代表 BFCL 任务比其他来源"更危险"——禁止关系多只是工具空间大的自然结果。

\subsection{确定性规则与人工复核}

授权标签由确定性规则生成，80 条抽样人工复核中 78 条与规则标签一致。2 条不一致来自语义边界情况：值匹配规则在"字段值恰好出现在参数中但语义上不该授权"的边界上存在局限。两名独立复核者的完整结果尚未提交，目前的人工一致性结论基于单侧复核。完整的人工双审结果出来后，可以进一步评估规则标签的可靠性。

\subsection{领域与伦理}

数据覆盖四个公开来源与多个应用领域。项目不主动收集现实个人数据；确定性补全值为虚构内容。数据适合研究用途，不应直接用作现实用户授权政策。AppWorld 衍生内容定位为私下科研交付；若要公开发布，须移除相应内容或采用符合上游条款的加密分发流程。
